\section{Dataset Extraction}
\label{sec:dataset}

\paragraph{Source.} GitSkills is distributed as a single SQLite file with four tables:
\texttt{artifacts}, with one row per file occurrence; \texttt{repos}, with repository
metadata; \texttt{artifact\_\allowbreak siblings}, with the files bundled with each skill; and
\texttt{mining\_runs}, with collection provenance~\cite{destefanis2027gitskills}. For this phase we use the official
sample released by the dataset authors~\cite{gitskills_sample}: the \ResContentsTotal{}
distinct contents with the lowest content hashes, together with every row the full dataset
holds about them, i.e., \ResOccTotal{} file occurrences in \ResReposTotal{} repositories owned
by \ResOwnersTotal{} accounts. Because content hashes are effectively uniform, the sample is a
random sample of distinct skills, and because every content group is complete, its copy counts
are exact. Queries written for the sample run unchanged on the full 41~GB file, which we will
use in Phase~3.

\paragraph{Fields.} Table~\ref{tab:fields} lists the fields we extract.

\begin{table}[t]
\caption{GitSkills fields used in this study.}
\label{tab:fields}
\small
\begin{tabular}{>{\raggedright\arraybackslash}p{0.55\columnwidth}>{\raggedright\arraybackslash}p{0.35\columnwidth}}
\toprule
Field(s) & Use \\
\midrule
\texttt{file\_sha}, \texttt{repo\_full\_name}, \texttt{path}, \texttt{location\_class} & occurrences, exact-copy counts \\
\texttt{filename}, \texttt{dedup\_primary}, \texttt{frontmatter\_valid}, \texttt{body\_chars} & scope filtering \\
\texttt{content}, \texttt{name}, \texttt{description} & skill text (representatives) \\
\texttt{first\_commit\_at} & ordering derived pairs \\
\texttt{artifact\_siblings.entry\_name} & bundled-file similarity \\
owner (prefix of \texttt{repo\_full\_name}) & same-owner and cross-owner reuse \\
\bottomrule
\end{tabular}
\end{table}

\paragraph{Scope.} The filename search behind GitSkills also returns lowercase
\texttt{skill.md} files and similar names that predate the format, so we keep only
representatives whose exact basename is \texttt{SKILL.md}. We require valid front matter
because the name and description are inputs to our classifier, and a body of at least 200
characters, which removes stubs and symbolic-link skills whose content is only a link target.
Table~\ref{tab:funnel} shows the effect of each filter.

\begin{table}[t]
\caption{Scope filters applied to the official GitSkills sample.}
\label{tab:funnel}
\small
\begin{tabular}{lr}
\toprule
Filter & Distinct skills \\
\midrule
All distinct contents & \ResFunnelZero \\
Exact basename \texttt{SKILL.md} & \ResFunnelOne \\
Valid YAML front matter & \ResFunnelTwo \\
Body $\geq$ 200 characters (dataset field) & \ResFunnelThree \\
Body $\geq$ 200 characters after our own parsing & \ResFunnelFour \\
\bottomrule
\end{tabular}
\end{table}

The resulting working set has \ResWorkContents{} distinct skills with \ResWorkOcc{} file
occurrences in \ResWorkRepos{} repositories owned by \ResWorkOwners{} accounts. Commit history
is available for \ResWorkHistory{} of these skills, and \ResWorkSiblings{} bundle at least one
other file. Already at this scale, \ResNamesTwoPlus{} front-matter names are carried by two or
more distinct contents.

\paragraph{Decomposition.} For each representative we split the file into its name
(normalised), description, prose (the body with fenced code removed), code blocks, and
heading sequence, and attach repository context: the owners of all its occurrences, the name
of the skill's folder, and the names of its bundled files. The result is stored as a small
SQLite database with two tables, \texttt{skills} and \texttt{occurrences}, which the remaining
steps read.
